\documentclass[11pt,a4paper]{article}
\usepackage{gtu-el-yazisi}
\begin{document}
\renewcommand\thempfootnote{\alph{mpfootnote}}% kutu içindeki dipnotlar için
\notbilgi{Diferansiyel Denklemler}{1. Hafta Uygulama (2. ders) -- Ayrılabilir ve homojen denklemler}{diff-2.ders-1-hafta.pdf}

\sayfa{1}
\hfill\fbox{1. HAFTA}
\section{Hatırlatma}
\begin{itemize}[label=\#]
\item $x$ bağımsız değişken, $y = y(x)$ ona bağlı fonksiyon ve bu fonksiyonun $y', y'', \dots, y^{(n)}$ türevleri arasındaki
\[
F\left(x, y, y', y'', \dots, y^{(n)}\right) = 0
\]
şeklindeki ilişkiye ``$n$. mertebeden adi diferansiyel denklem'' denir.
\item Eğer $F$ fonksiyonu $y$ ve onun $y', y'', \dots, y^{(n)}$ türevlerine göre lineer olursa bu durumda denkleme ``lineer adi diferansiyel denklem'' denir.
\item \underline{Ayrılabilir denklemler}: \ $y' = f_1(x)\cdot f_2(y)$ şeklinde veya
$M_1(x)\cdot M_2(y)\,dx + N_1(x)\cdot N_2(y)\,dy = 0$ şeklindeki denklemlere ``değişkenlerine ayrılabilir denklem'' denir.

Bu durumda, $M_2(y)\cdot N_1(x) \neq 0$ olursa, denklemi bu çarpana bölüp, integral alınırsa diferansiyel denklemin genel çözümü elde edilmiş olur. Tabi, bu durumda $M_2(y)\cdot N_1(x) = 0$ durumu için çözüm kaybolabilir. O yüzden bu durumu ayrı incelemek gerekir.
\item \underline{Ayrılabilir denklemlere indirgenebilir denklemler}: iki tip denklem vardır.
\begin{itemize}[label={}]
\item (a) $y' = f(ax+by+c) \Rightarrow z = ax+by+c$ dönüşümü yapılarak integral alınır.
\item (b) Homojen denklemler: \ $y' = f\!\left(\dfrac{y}{x}\right) \Rightarrow u = \dfrac{y}{x}$ dönüşümü yapılarak integral alınır.
\end{itemize}
\end{itemize}

\section{Alıştırmalar}
\begin{ornek}[1]
Aşağıda verilen diferansiyel denklemlerin türlerini mertebelerine göre belirleyiniz
\end{ornek}
\begin{itemize}[label=$\ast$]
\item $\dfrac{dy}{dx} + x^2y = xe^{2}$ \qquad $\rightarrow$ 1. mertebeden lineer denklem
\item $\dfrac{d^4y}{dx^4} + \left(\dfrac{d^2y}{dx^2}\right)^5 + 5y = 0$ \quad $\rightarrow$ 4. mertebeden lineer olmayan denklem
\end{itemize}

\sayfa{2}
\begin{itemize}[label=$\ast$]
\item $\dfrac{d^2y}{dx^2} + y\sin x = 0$ \qquad $\rightarrow$ 2. mertebeden lineer denklem
\item $\dfrac{d^6x}{dt^6} + \left(\dfrac{d^4x}{dt^4}\right)\left(\dfrac{d^3x}{dt^3}\right) + x = t$ \quad $\rightarrow$ 6. mertebeden lineer olmayan denklem
\item $\dfrac{d^3y}{dx^3} + 4\dfrac{d^2y}{dx^2} - 5\dfrac{dy}{dx} + 3y = \sin x$ \quad $\rightarrow$ 3. mertebeden lineer denklem
\end{itemize}

\begin{ornek}[2]
$x^3 + 3xy^2 = 1$ fonksiyonunun $0<x<1$ aralığında $2xy\left(\dfrac{dy}{dx}\right) + x^2 + y^2 = 0$ diferansiyel denkleminin bir kapalı çözümü olduğunu gösteriniz.
\end{ornek}
\begin{cozum}
$x^3 + 3xy^2 = 1$ ifadesinin $x$ e göre türevini alalım.
\[
3x^2 + 3y^2 + 6xyy' = 0 \Rightarrow 3\left(x^2 + y^2 + 2xyy'\right) = 0 \Rightarrow \kutula{x^2 + y^2 + 2xyy' = 0}
\]
sağlandığından verilen kapalı form diferansiyel denklemin bir çözümüdür.
\end{cozum}

\begin{ornek}[3]
$f(x) = e^{mx}$ fonksiyonunun $\dfrac{d^3y}{dx^3} - 3\dfrac{d^2y}{dx^2} - 4\dfrac{dy}{dx} + 12y = 0$ diferansiyel denkleminin bir çözümü olması için $m$ nin alacağı değerler nedir?
\end{ornek}
\begin{cozum}
$y = e^{mx}$, \ $y' = me^{mx}$, \ $y'' = m^2e^{mx}$, \ $y''' = m^3e^{mx}$ diferansiyel denklemde yerine koyarsak,
\[
m^3e^{mx} - 3m^2e^{mx} - 4me^{mx} + 12e^{mx} = 0 \Rightarrow \overbrace{e^{mx}}^{\neq 0}\left(m^3 - 3m^2 - 4m + 12\right) = 0
\]
\[
\Rightarrow m^3 - 3m^2 - 4m + 12 = 0 \Rightarrow m^2(m-3) - 4(m-3) = 0 \Rightarrow (m-3)(m-2)(m+2) = 0
\]
\[
\Rightarrow m = \pm2,\ 3
\]
elde edilir. O halde diferansiyel denklemin çözümleri $y = e^{2x}, e^{-2x}, e^{3x}$ şeklindedir.
\end{cozum}

\begin{ornek}[4]
$5x^2y^2 - 2x^3y^2 = 1$ fonksiyonunun $x\left(\dfrac{dy}{dx}\right) + y = x^3y^3$ fonksiyonunun $0<x<\dfrac52$ aralığında denklemin çözümü olduğunu gösteriniz.
\end{ornek}

\sayfa{3}
\begin{cozum}
$5x^2y^2 - 2x^3y^2 = 1$ fonksiyonunun $x$ e göre türevini alalım.
\[
10xy^2 + 10x^2yy' - 6x^2y^2 - 4x^3yy' = 0 \Rightarrow xy'\left(10xy - 4x^2y\right) = 6x^2y^2 - 10xy^2
\]
\[
\Rightarrow xy' = \frac{2xy\left(3xy - 5y\right)}{2xy\left(5 - 2x\right)} \Rightarrow xy' = \frac{3xy - 5y}{5 - 2x}
\]
ifadesini diferansiyel denklemde yerine yazalım:
\[
\left.\begin{aligned}
&xy' + y = \frac{3xy - 5y}{5-2x} + y = \frac{3xy - 5y + 5y - 2xy}{5-2x} = \frac{xy}{5-2x}\\
&5x^2y^2 - 2x^3y^2 = 1 \Rightarrow x^2y^2(5 - 2x) = 1 \Rightarrow 5 - 2x = \frac{1}{x^2y^2}
\end{aligned}\right\} \Rightarrow
\]
\[
xy' + y = \frac{xy}{\frac{1}{x^2y^2}} \Rightarrow \kutula{xy' + y = x^3y^3}
\]
elde edilir. Böylece çözüm olduğu gösterilmiş olur.
\end{cozum}

\begin{ornek}[5]
$f(x) = \dfrac{1}{1+x^2}$ nin $(1+x^2)\dfrac{d^2y}{dx^2} + 4x\dfrac{dy}{dx} + 2y = 0$ diferansiyel denkleminin çözümü olduğunu gösteriniz.
\end{ornek}
\begin{cozum}
\[
y = \frac{1}{1+x^2}, \qquad y' = \frac{-2x}{(1+x^2)^2}, \qquad y'' = \frac{-2(1+x^2)^2 + 8x^2(1+x^2)}{(1+x^2)^4} = \frac{6x^2-2}{(1+x^2)^3}
\]
ifadelerini diferansiyel denklemde yerine koyarsak,
\[
\frac{(1+x^2)(6x^2-2)}{(1+x^2)^3} + 4x\cdot\left(\frac{-2x}{(1+x^2)^2}\right) + \underset{(1+x^2)}{\frac{2}{1+x^2}} = \frac{6x^2 - 2 - 8x^2 + 2 + 2x^2}{(1+x^2)^2} = 0
\]
olduğundan verilen $f(x)$ fonksiyonu diferansiyel denklemin bir çözümüdür.
\end{cozum}

\begin{ornek}[6]
Aşağıda verilen diferansiyel denklemleri çözünüz.

\begin{tabular}{@{}ll@{}}
(a) $4xy\,dx + (x^2+1)\,dy = 0$ & \qquad (c) $y' + y^2\sin x = 0$ \\
(b) $\operatorname{cosec}y\,dx + \sec x\,dy = 0$ & \qquad (d) $y' = (\cos^2x)(\cos^2 2y)$
\end{tabular}
\end{ornek}

\sayfa{4}
\begin{cozum}
\textbf{(a)} \ $y(x^2+1) \neq 0$ olmak üzere,
\[
\frac{4x}{x^2+1}dx + \frac1y dy = 0 \Rightarrow 2\int\frac{2x\,dx}{x^2+1} + \int\frac1y dy = c_1 \underset{\substack{\downarrow\\ x^2+1=u\\ 2x\,dx=du}}{\Rightarrow} 2\ln(x^2+1) + \ln|y| = c_1
\]
\[
\Rightarrow \ln(x^2+1)^2|y| = c_1 \Rightarrow |y| = \frac{e^{c_1}}{(x^2+1)^2}
\]
\[
\underset{e^{c_1}=c>0}{\Rightarrow} y = \frac{c}{(x^2+1)^2}, \quad c \in \mathbb{R}
\]
Şimdi, kaybolma ihtimali olan çözümlere bakalım:

$y(x^2+1) = 0 \Rightarrow y = 0 \rightarrow$ Bu çözümü yukarıda $c = 0$ yazarak elde edebiliriz.
O halde, çözüm ailesi
\[
y = \frac{c}{(x^2+1)^2}, \quad c \in \mathbb{R}
\]
şeklindedir.

\textbf{(b)} \ $\operatorname{cosec}y\cdot\sec x \neq 0$ olmak üzere
\[
\frac{dx}{\sec x} + \frac{dy}{\operatorname{cosec}y} = 0 \Rightarrow \cos x\,dx + \sin y\,dy = 0 \Rightarrow \int\cos x\,dx + \int\sin y\,dy = c
\]
\[
\Rightarrow \kutula{\sin x - \cos y = c}, \quad c \in \mathbb{R}
\]

\textbf{(c)} \ $\dfrac{dy}{dx} + y^2\sin x = 0 \Rightarrow dy + y^2\sin x\,dx = 0$ denklemini çözelim.

$y^2 \neq 0$ olmak üzere,
\[
\frac{dy}{y^2} + \sin x\,dx = 0 \Rightarrow \int\frac{dy}{y^2} + \int\sin x\,dx = -c \Rightarrow -\frac1y - \cos x = -c \Rightarrow \frac1y + \cos x = c
\]
Şimdi, kaybolma ihtimali olan çözümlere bakalım.

$y^2 = 0 \Rightarrow y = 0$ bir çözümdür.
\sayfa{5}
O halde, çözüm ailesi
\[
\begin{cases}
\dfrac1y + \cos x = c \ ; & c \in \mathbb{R}\\[2mm]
y = 0
\end{cases}
\]
şeklindedir.

\textbf{(d)} \ $\dfrac{dy}{dx} = (\cos^2x)(\cos^2 2y) \Rightarrow dy = (\cos^2x)(\cos^2 2y)\,dx$ denklemini çözelim.

$\cos^2 2y \neq 0$ olmak üzere,
\[
\frac{dy}{\cos^2 2y} - \cos^2x\,dx = 0 \Rightarrow \int\frac{dy}{\cos^2 2y} - \int\cos^2x\,dx = c_1
\]
\[
\Rightarrow \int\frac{dy}{\cos^2 2y} - \int\frac{\cos 2x + 1}{2}\,dx = c_1 \Rightarrow \frac{\tan 2y}{2} - \frac{\sin 2x}{4} - \frac{x}{2} = c_1
\]
\[
\underset{\substack{\downarrow\\ (2c_1=c)}}{\Rightarrow} 2\tan 2y - \sin 2x - 2x = c, \quad c \in \mathbb{R} \qquad\text{\footnotemark}
\]
\footnotetext{Claude notu: denklem 4 ile çarpıldığı için muhtemelen $4c_1 = c$ olmalı.}
Şimdi, kaybolma ihtimali olan çözümlere bakalım.
\[
\cos^2 2y = 0 \Rightarrow 2y = \pm\frac{(2n+1)\pi}{2} \Rightarrow y = \pm\frac{(2n+1)\pi}{4}, \quad \forall n \in \mathbb{N} \rightarrow \text{çözümdür}
\]
O halde, çözüm ailesi,
\[
\begin{cases}
2\tan 2y - \sin 2x - 2x = c, & c \in \mathbb{R}\\[1mm]
y = \pm\dfrac{(2n+1)}{4}\pi, & n \in \mathbb{N}
\end{cases}
\]
şeklindedir.
\end{cozum}

\sayfa{6}
\begin{ornek}[7]
Aşağıdaki başlangıç-değer (Cauchy) problemlerini çözünüz.

(a) $(y+2)\,dx + y(x+4)\,dy = 0$, \quad $y(-3) = -1$

(b) $y' = \dfrac{x(x^2+1)}{4y^3}$, \quad $y(0) = -\dfrac{1}{\sqrt2}$

(c) $(3x+8)(y^2+4)\,dx - 4y(x^2+5x+6)\,dy = 0$, \quad $y(1) = 2$
\end{ornek}
\begin{cozum}
\textbf{(a)} \ $(y+2)(x+4) \neq 0$ olmak üzere,
\[
\frac{dx}{x+4} + \frac{y}{y+2}dy = 0 \Rightarrow \int\frac{dx}{x+4} + \int\frac{\overset{\nearrow\, y=y+2-2}{y}}{y+2}dy = c
\]
\[
\Rightarrow \int\frac{dx}{x+4} + \int\left(1 - \frac{2}{y+2}\right)dy = c \Rightarrow \ln|x+4| + y - 2\ln|y+2| = c
\]
\[
\Rightarrow \kutula{y + \ln\frac{|x+4|}{(y+2)^2} = c}, \quad c \in \mathbb{R}
\]
Şimdi, kaybolma ihtimali olan çözümlere bakalım.

$(y+2)(x+4) = 0 \Rightarrow y = -2 \rightarrow$ çözümdür.

O halde, çözüm ailesi
\[
\begin{cases}
y + \ln\dfrac{|x+4|}{(y+2)^2} = c\\[2mm]
y = -2, \ x = -4
\end{cases}
\]
şeklindedir. Şimdi, $y(-3) = -1$ koşulunu sağlayan çözümü bulalım.
\[
-1 + \ln\frac{1}{(-1+2)^2} = c \Rightarrow -1 = c
\]
\[
\Rightarrow y + \ln\frac{|x+4|}{(y+2)^2} = -1 \Rightarrow \kutula{e^{-(y+1)}(y+2)^2 = |x+4|} \ /\!/
\]
\sayfa{7}
\textbf{(b)} \ $\dfrac{dy}{dx} = \dfrac{x(x^2+1)}{4y^3} \Rightarrow 4y^3dy - x(x^2+1)\,dx = 0$
\[
\Rightarrow 4\int y^3dy - \int x(x^2+1)\,dx = c_1 \underset{\substack{\downarrow\\ x^2+1=u\\ 2x\,dx=du}}{\Rightarrow} 4\frac{y^4}{4} - \frac{(x^2+1)^2}{4} = c_1 \overset{\substack{(4c_1=c)\\ \uparrow}}{\Rightarrow} \kutula{4y^4 - (x^2+1)^2 = c}
\]
Şimdi, $y(0) = -\dfrac{1}{\sqrt2}$ koşulunu sağlayan çözümü bulalım.
\[
4\cdot\left(-\frac{1}{\sqrt2}\right)^4 - (0+1)^2 = c \Rightarrow c = 0 \Rightarrow 4y^4 - (x^2+1)^2 = 0 \Rightarrow \kutula{y = \pm\sqrt{\frac{x^2+1}{2}}} \ /\!/
\]

\textbf{(c)} \ $(y^2+4)(x^2+5x+6) \neq 0$ olmak üzere,
\[
\frac{3x+8}{x^2+5x+6}dx - \frac{4y}{y^2+4}dy = 0 \underset{\substack{\downarrow\\ \left(\frac{3x+8}{x^2+5x+6} = \frac{1}{x+3} + \frac{2}{x+2}\right)}}{\Rightarrow} \frac{dx}{x+3} + \frac{2\,dx}{x+2} - \frac{4y}{y^2+4}dy = 0
\]
\[
\Rightarrow \int\frac{dx}{x+3} + 2\int\frac{dx}{x+2} - 2\int\frac{2y\,dy}{y^2+4} = c_1 \Rightarrow \ln|x+3| + 2\ln|x+2| - 2\ln|y^2+4| = c_1
\]
\[
\Rightarrow \ln\frac{|x+3|(x+2)^2}{(y^2+4)^2} = c_1 \underset{\substack{\downarrow\\ (e^{c_1}=c)}}{\Rightarrow} \kutula{c\,(y^2+4)^2 = |x+3|(x+2)^2 \quad c \in \mathbb{R}}
\]
Şimdi, $y(1) = 2$ koşulunu sağlayan çözümü bulalım.
\[
c\,8^2 = 4\cdot9 \Rightarrow c = \frac{9}{16} \Rightarrow \kutula{(x+3)(x+2)^2 = \frac{9}{16}(y^2+4)^2} \ /\!/
\]
\end{cozum}

\sayfa{8}
\begin{ornek}[8]
Aşağıdaki diferansiyel denklemleri çözünüz.

(a) $(2xy + 3y^2)\,dx - (2xy + x^2)\,dy = 0$

(b) $\left(x^3 + y^2\sqrt{x^2+y^2}\right)dx - \left(xy\sqrt{x^2+y^2}\right)dy = 0$

(c) $(x^2 + 3y^2)\,dx - 2xy\,dy = 0$, \quad $y(2) = 6$
\end{ornek}
\begin{cozum}
$M(x,y) = 2xy + 3y^2$ ve $N(x,y) = 2xy + x^2$ olmak üzere,
\[
\left.\begin{aligned}
M(tx,ty) &= 2\cdot tx\cdot ty + 3\cdot t^2y^2 = t^2\left(2xy + 3y^2\right) = t^2M(x,y)\\
N(tx,ty) &= 2\cdot tx\cdot ty + t^2x^2 = t^2\left(2xy + x^2\right) = t^2N(x,y)
\end{aligned}\right\} \Rightarrow \begin{array}{l}\text{2. dereceden}\\ \text{homojen denklem}\end{array}
\]
Şimdi, $\dfrac{y}{x} = u$ dönüşümü yapalım.
\[
\frac{y}{x} = u \Rightarrow y = ux \Rightarrow y' = u'x + u \Rightarrow \frac{dy}{dx} = \frac{du}{dx}x + u \Rightarrow dy = du\,x + u\,dx
\]
\[
\Rightarrow \left(2x\cdot ux + 3u^2x^2\right)dx - \left(2xux + x^2\right)\left(x\,du + u\,dx\right) = 0
\]
\[
\Rightarrow \left(2x^2u + 3u^2x^2 - 2x^2u^2 - x^2u\right)dx + \left(-2x^3u - x^3\right)du = 0
\]
\[
\Rightarrow \left(x^2u^2 + x^2u\right)dx + \left(-2x^3u - x^3\right)du \Rightarrow x^2\left(u^2+u\right)dx - x^3(2u+1)\,du = 0
\]
$(u^2+u)x^3 \neq 0$ olmak üzere,
\[
\frac1x dx - \frac{2u+1}{u^2+u}du = 0 \Rightarrow \int\frac1x dx - \int\frac{2u+1}{u^2+u}du = c_1 \Rightarrow \ln|x| - \ln|u^2+u| = c_1
\]
\[
\underset{e^{c_1}=c_2}{\Rightarrow} \left|\frac{x}{u^2+u}\right| = c_2 \Rightarrow \frac{x}{u^2+u} = \overbrace{(\pm c_2)}^{c} \Rightarrow x = c\left(u^2+u\right)
\]
\[
\underset{}{\overset{\frac{y}{x}=u}{\Rightarrow}} x = c\left(\frac{y^2}{x^2} + \frac{y}{x}\right) \qquad c \in \mathbb{R}
\]
Kaybolan çözümler: \ $u^2 + u = 0 \Rightarrow u(u+1) = 0 \Rightarrow
\left.\begin{aligned}&u = 0 \Rightarrow y = 0\\ &u = -1 \Rightarrow y = -x\end{aligned}\right\}$ çözümdür.
\sayfa{9}
\textbf{(b)} \ $M(x,y) = \left(x^3 + y^2\sqrt{x^2+y^2}\right)$ ve $N(x,y) = xy\sqrt{x^2+y^2}$ olmak üzere,
\[
\left.\begin{aligned}
M(tx,ty) &= t^3x^3 + t^2y^2\sqrt{t^2(x^2+y^2)} = t^3M(x,y)\\
N(tx,ty) &= tx\cdot ty\sqrt{t^2(x^2+y^2)} = t^3N(x,y)
\end{aligned}\right\} \Rightarrow \text{3. dereceden homojen denklem}
\]
Şimdi, $\dfrac{y}{x} = u$ dönüşümü yapalım.
\[
y = ux \Rightarrow y' = u'x + u \Rightarrow \frac{dy}{dx} = \frac{du}{dx}x + u \Rightarrow dy = x\,du + u\,dx
\]
\[
\left(x^3 + u^2x^2\sqrt{x^2+u^2x^2}\right)dx - \left(xux\sqrt{x^2+u^2x^2}\right)\left(x\,du + u\,dx\right) = 0
\]
\[
\Rightarrow \left(x^3 + u^2x^2\sqrt{x^2+u^2x^2} - x^2u^2\sqrt{x^2+u^2x^2}\right)dx - x^3u\sqrt{x^2+u^2x^2}\,du = 0
\]
\[
\Rightarrow x^3dx - ux^4\sqrt{1+u^2}\,du = 0
\]
$x^4 \neq 0$ olmak üzere,
\[
\frac1x dx - u\sqrt{1+u^2}\,du = 0 \Rightarrow \int\frac1x dx - \int u\sqrt{1+u^2}\,du = c_1
\]
\[
\Rightarrow \ln|x| - \frac{(1+u^2)^{3/2}}{3} = c_1 \underset{\frac{y}{x}=u}{\Rightarrow} \ln|x| = c_1 + \frac{\left(1+\left(\frac{y}{x}\right)^2\right)^{3/2}}{3}
\]
\[
\Rightarrow 3\ln|x| = 3c_1 + \frac{(x^2+y^2)^{3/2}}{(x^2)^{3/2}} \underset{\substack{\downarrow\\ 3c_1=c}}{\Rightarrow} \kutula{x^3\ln|x^3| = c + (x^2+y^2)^{3/2}} \quad\text{\footnotemark}
\]
\footnotetext{Claude notu: muhtemelen $x^3\ln|x^3| = c\,x^3 + (x^2+y^2)^{3/2}$ olmalı ($c$ terimi de $x^3$ ile çarpılmalı).}

\textbf{(c)} \ $M(x,y) = x^2 + 3y^2$ ve $N(x,y) = 2xy$ olmak üzere,
\[
\left.\begin{aligned}
M(tx,ty) &= t^2M(x,y)\\
N(tx,ty) &= t^2N(x,y)
\end{aligned}\right\} \Rightarrow \text{2. dereceden homojen denklem}
\]
Şimdi, $\dfrac{y}{x} = u$ dönüşümü yapalım.
\[
y = ux \Rightarrow dy = x\,du + u\,dx
\]
\sayfa{10}
\[
\left(x^2 + 3u^2x^2\right)dx - (2xux)\left(u\,dx + x\,du\right) = 0
\]
\[
\Rightarrow \left(x^2 + 3u^2x^2 - 2x^2u^2\right)dx - 2x^3u\,du = 0
\]
\[
\Rightarrow \left(x^2 + u^2x^2\right)dx - 2x^3u\,du = 0 \Rightarrow x^2\left(1+u^2\right)dx - 2x^3u\,du = 0
\]
$(1+u^2)x^3 \neq 0$ olmak üzere,
\[
\frac1x dx - \frac{2u}{1+u^2}du = c \quad\text{\footnotemark} \Rightarrow \int\frac1x dx - \int\frac{2u}{1+u^2}du = c_1 \Rightarrow \ln|x| - \ln(1+u^2) = c_1
\]
\footnotetext{Claude notu: bu adımda sağ taraf muhtemelen $0$ olmalı (integral alındıktan sonra sabit $c_1$ geliyor).}
\[
\underset{e^{c_1}=c}{\Rightarrow} \frac{x}{1+u^2} = c \Rightarrow x = c\left(1+u^2\right) \underset{u=\frac{y}{x}}{\Rightarrow} \kutula{x^3 = c\left(x^2+y^2\right) \quad c \in \mathbb{R}}
\]
Şimdi, $y(2) = 6$ koşulunu sağlayan çözümü bulalım.
\[
8 = c\,(4+36) \Rightarrow c = \frac15 \Rightarrow \kutula{x^3 = \frac15\left(x^2+y^2\right)} \ /\!/
\]
\end{cozum}

\end{document}
